\RequirePackage{plautopatch}
\documentclass[platex,dvipdfmx,a4paper]{jsarticle}

\usepackage{amsmath}
\usepackage{amsfonts}
\usepackage{amssymb}

\usepackage[dvipdfmx]{graphicx}
\usepackage{float}

\usepackage{tikz}
\usepackage{caption}

\usepackage[%
  dvipdfmx,%
  setpagesize=false,%
  bookmarks=true,%
  bookmarksdepth=tocdepth,%
  bookmarksnumbered=true,%
  colorlinks=false,%
  pdftitle={},%
  pdfsubject={},%
  pdfauthor={},%
  pdfkeywords={}%
]{hyperref}
\usepackage{pxjahyper}

\usepackage{framed}
\setlength{\OuterFrameSep}{0pt}

\usepackage{physics2}
\usephysicsmodule{ab}
\usephysicsmodule{ab.braket}

\newcommand{\hb}[2]{\frac{\partial #1}{\partial #2}}
\newcommand{\bbn}[2]{\frac{d #1}{d #2}}
\newcommand{\nhb}[3]{\frac{\partial^{#1} #2}{\partial #3^{#1}}}
\newcommand{\nbbn}[3]{\frac{d^{#1} #2}{d #3^{#1}}}
\newcommand{\vectxy}[2]{\begin{bmatrix}#1\#2\end{bmatrix}}
\newcommand{\vectxyz}[3]{\begin{bmatrix}#1\#2\#3\end{bmatrix}}
\newcommand{\vect}[1]{\boldsymbol{\mathrm{#1}}}

\newcommand{\ssum}[3]{\sum_{#1=#2}^#3}

\newcommand{\veff}{V_{\mathrm{eff}}(r)}

\hypersetup{colorlinks=true}
\hypersetup{linkcolor=red}

\title{量子力学入門演習レポート5}
\author{25B00729　嶋田 悠音}

\begin{document}
\maketitle

\subsection*{7.2}
\subsection*{(1)}

まず、任意の演算子 \(\hat{X}\)、\(\hat{Y}\)、\(\hat{Z}\) に対する交換子の積の公式
\begin{align*}
  [\hat{X},\hat{Y}\hat{Z}]=[\hat{X},\hat{Y}]\hat{Z}+\hat{Y}[\hat{X},\hat{Z}]
\end{align*}
を用いて、
\begin{align*}
  [\hat{a},(\hat{a}^{\dagger})^n]=n(\hat{a}^{\dagger})^{n-1}
\end{align*}
を数学的帰納法によって示す。\(n=1\) の場合は、与えられた交換関係から
\begin{align*}
  [\hat{a},\hat{a}^{\dagger}]=1
\end{align*}
であるため、主張は成り立つ。次に、ある正の整数 \(n\) に対して
\begin{align*}
  [\hat{a},(\hat{a}^{\dagger})^n]=n(\hat{a}^{\dagger})^{n-1}
\end{align*}
が成り立つと仮定する。このとき、交換子の積の公式を用いる：
\begin{align*}
  [\hat{a},(\hat{a}^{\dagger})^{n+1}]&=[\hat{a},(\hat{a}^{\dagger})^n]\hat{a}^{\dagger}+(\hat{a}^{\dagger})^n[\hat{a},\hat{a}^{\dagger}]=n(\hat{a}^{\dagger})^n+(\hat{a}^{\dagger})^n=(n+1)(\hat{a}^{\dagger})^n。
\end{align*}
したがって、数学的帰納法により、すべての正の整数 \(n\) に対して
\begin{align*}
  [\hat{a},(\hat{a}^{\dagger})^n]=n(\hat{a}^{\dagger})^{n-1}
\end{align*}
が成り立つ。
次に、交換子の積の公式と先ほど示した結果を用いる：
\begin{align*}
  [\hat{a}^{\dagger}\hat{a},(\hat{a}^{\dagger})^n]&=\hat{a}^{\dagger}[\hat{a},(\hat{a}^{\dagger})^n]+[\hat{a}^{\dagger},(\hat{a}^{\dagger})^n]\hat{a}=n(\hat{a}^{\dagger})^n+0=n(\hat{a}^{\dagger})^n。
\end{align*}
よって、
\begin{align*}
  [\hat{a}^{\dagger}\hat{a},(\hat{a}^{\dagger})^n]=n(\hat{a}^{\dagger})^n
\end{align*}
を得る。
また、最初に示した交換関係の Hermite 共役を取ると、
\begin{align*}
  [\hat{a}^n,\hat{a}^{\dagger}]=n\hat{a}^{n-1}
\end{align*}
が成り立つ。したがって、
\begin{align*}
  [\hat{a}^{\dagger},\hat{a}^n]=-n\hat{a}^{n-1}
\end{align*}
である。これと交換子の積の公式を用いる：
\begin{align*}
  [\hat{a}^{\dagger}\hat{a},\hat{a}^n]&=\hat{a}^{\dagger}[\hat{a},\hat{a}^n]+[\hat{a}^{\dagger},\hat{a}^n]\hat{a}=-n\hat{a}^n。
\end{align*}
よって、
\begin{align*}
  [\hat{a}^{\dagger}\hat{a},\hat{a}^n]=-n\hat{a}^n
\end{align*}
が示された。なお、\(n=0\) の場合はいずれも恒等演算子との交換子となるため、両辺は \(0\) である。

\subsection*{(2)}

状態 \(\ket|0>\) および \(\ket|n>\) は、
\begin{align*}
  \hat{a}\ket|0>=0,\qquad \braket<0|0>=1,\qquad \ket|n>=C_n(\hat{a}^{\dagger})^n\ket|0>
\end{align*}
によって定義される。まず、\(\hat{a}\ket|0>=0\) より、
\begin{align*}
  \hat{N}\ket|0>=\hat{a}^{\dagger}\hat{a}\ket|0>=0
\end{align*}
が成り立つ。交換子の定義から、
\begin{align*}
  \hat{N}(\hat{a}^{\dagger})^n=(\hat{a}^{\dagger})^n\hat{N}+[\hat{N},(\hat{a}^{\dagger})^n]=(\hat{a}^{\dagger})^n\hat{N}+n(\hat{a}^{\dagger})^n
\end{align*}
である。したがって、\(\ket|n>\) に \(\hat{N}\) を作用させる：
\begin{align*}
  \hat{N}\ket|n>&=C_n\hat{N}(\hat{a}^{\dagger})^n\ket|0>=C_n\left\{(\hat{a}^{\dagger})^n\hat{N}+n(\hat{a}^{\dagger})^n\right\}\ket|0>=nC_n(\hat{a}^{\dagger})^n\ket|0>=n\ket|n>。
\end{align*}
よって、状態 \(\ket|n>\) は数演算子 \(\hat{N}=\hat{a}^{\dagger}\hat{a}\) の固有状態であり、その固有値は \(n\) である。

\subsection*{(3)}

規格化条件 \(\braket<n|n>=1\) から定数 \(C_n\) を定める。\(\ket|n>=C_n(\hat{a}^{\dagger})^n\ket|0>\) の Hermite 共役を取ると、
\begin{align*}
  \bra<n|=C_n^*\bra<0|\hat{a}^n
\end{align*}
であるため、
\begin{align*}
  \braket<n|n>=|C_n|^2\braket<0|\hat{a}^n(\hat{a}^{\dagger})^n|0>
\end{align*}
を得る。ここで、(1) の結果から、
\begin{align*}
  \hat{a}(\hat{a}^{\dagger})^n=(\hat{a}^{\dagger})^n\hat{a}+n(\hat{a}^{\dagger})^{n-1}
\end{align*}
である。最右辺の \(\hat{a}\) が \(\ket|0>\) に作用する項は消えるため、
\begin{align*}
  \braket<0|\hat{a}^n(\hat{a}^{\dagger})^n|0>&=n\braket<0|\hat{a}^{n-1}(\hat{a}^{\dagger})^{n-1}|0>=n(n-1)\braket<0|\hat{a}^{n-2}(\hat{a}^{\dagger})^{n-2}|0>=\cdots=n!\braket<0|0>=n!。
\end{align*}
したがって、規格化条件は
\begin{align*}
  1=|C_n|^2n!
\end{align*}
となるので、
\begin{align*}
  |C_n|=\frac{1}{\sqrt{n!}}
\end{align*}
を得る。\(C_n\) を正の実数に選ぶと、
\begin{align*}
  C_n=\frac{1}{\sqrt{n!}}。
\end{align*}
よって、規格化された状態は
\begin{align*}
  \ket|n>=\frac{1}{\sqrt{n!}}(\hat{a}^{\dagger})^n\ket|0>
\end{align*}
と表される。

\subsection*{(4)}

まず、生成演算子 \(\hat{a}^{\dagger}\) を \(\ket|n>\) に作用させる：
\begin{align*}
  \hat{a}^{\dagger}\ket|n>&=\frac{1}{\sqrt{n!}}(\hat{a}^{\dagger})^{n+1}\ket|0>=\sqrt{n+1}\frac{1}{\sqrt{(n+1)!}}(\hat{a}^{\dagger})^{n+1}\ket|0>=\sqrt{n+1}\ket|n+1>。
\end{align*}
したがって、
\begin{align*}
  A_n=\sqrt{n+1}
\end{align*}
である。

次に、\(n\geq1\) に対して消滅演算子 \(\hat{a}\) を \(\ket|n>\) に作用させる。交換子の定義と (1) の結果から、
\begin{align*}
  \hat{a}(\hat{a}^{\dagger})^n=(\hat{a}^{\dagger})^n\hat{a}+n(\hat{a}^{\dagger})^{n-1}
\end{align*}
であるため、
\begin{align*}
  \hat{a}\ket|n>&=\frac{1}{\sqrt{n!}}\hat{a}(\hat{a}^{\dagger})^n\ket|0>=\frac{n}{\sqrt{n!}}(\hat{a}^{\dagger})^{n-1}\ket|0>=\sqrt{n}\frac{1}{\sqrt{(n-1)!}}(\hat{a}^{\dagger})^{n-1}\ket|0>=\sqrt{n}\ket|n-1>。
\end{align*}
また、\(n=0\) の場合は定義より \(\hat{a}\ket|0>=0\) である。したがって、
\begin{align*}
  B_n=\sqrt{n}
\end{align*}
を得る。以上より、昇降演算子の作用は
\begin{align*}
  \hat{a}^{\dagger}\ket|n>=\sqrt{n+1}\ket|n+1>,\qquad \hat{a}\ket|n>=\sqrt{n}\ket|n-1>
\end{align*}
であり、\(\hat{a}^{\dagger}\) は数演算子の固有値を \(1\) 増加させ、\(\hat{a}\) は固有値を \(1\) 減少させることが示された。

\end{document}